\respcap{cap:sortcomp}
\resp{sc-merge} Ver o procedimento \textsc{Intercala} da §\ref{sec:sortcomp}: dois índices percorrem as metades, o menor elemento apontado vai para o vetor auxiliar (``$\le$'' para manter a estabilidade), o que sobra de uma das metades é copiado ao final e o auxiliar volta para o vetor. Custo $\Theta(n)$.

\resp{sc-iter} Intercale blocos de tamanho $1,2,4,\dots$:
\begin{pseudo}{Mergesort iterativo (de baixo para cima)}
\begin{algorithmic}
\State $tam:=1$
\While{$tam<n$}
  \State $ini:=1$
  \While{$ini+tam\le n$}
    \State $meio:=ini+tam-1$;\ $fim:=\min\{ini+2\,tam-1,\ n\}$
    \State \Call{Intercala}{$V,ini,meio,fim$};\ $ini:=ini+2\,tam$
  \EndWhile
  \State $tam:=2\,tam$
\EndWhile
\end{algorithmic}
\end{pseudo}
São $\ceil{\log_2n}$ passadas de custo $O(n)$: $\Theta(n\log n)$, sem pilha de recursão.

\resp{sc-semcomp} Se as chaves são inteiros em $[0,k]$, conte as ocorrências de cada valor e use as contagens para posicionar os elementos (\emph{counting sort}, $\Theta(n+k)$). Se têm $d$ dígitos na base $b$, ordene dígito a dígito, do menos significativo para o mais significativo, com um método estável (\emph{radix sort}, $\Theta(d(n+b))$). Se são reais uniformemente distribuídos em $[0,1)$, distribua-os em $n$ baldes (\emph{bucket sort}, $\Theta(n)$ esperado). Todas usam a estrutura das chaves (Capítulo~\ref{cap:sortlin}); para chaves arbitrárias, só resta comparar.

\resp{sc-estavel} (a) Sim: na intercalação, em caso de empate, copie o elemento da metade esquerda (teste ``$\le$''). Por indução, as metades estão ordenadas de forma estável, e a intercalação preserva a ordem relativa dos iguais, que vêm todos da esquerda antes dos da direita. (b) A versão in-place usual não é estável: em $[2_a,2_b,1]$, com pivô 1, a troca de $2_a$ com 1 coloca $2_a$ depois de $2_b$. É possível torná-lo estável com uma partição que copia os menores, os iguais e os maiores para vetores auxiliares na ordem original, ao custo de $\Theta(n)$ de memória extra, ou usando como chave o par (valor, posição original).

\resp{sc-mestre} (a) $E=2$ e $f=n=O(n^{2-\varepsilon})$: $\Theta(n^2)$. (b) $f=\Theta(n^2)$: $\Theta(n^2\log n)$. (c) $f=\Omega(n^{2+\varepsilon})$ e $4(n/2)^3=n^3/2$: $\Theta(n^3)$. (d) $E=1/2$ e $f=\Theta(\sqrt n)$: $\Theta(\sqrt n\log n)$. (e) $E=2$: $\Theta(n^2)$. (f) $E=0$, $f=n$ e $n/2\le cn$: $\Theta(n)$.

\resp{sc-arvore} Cada nível custa no máximo $n$, pois os subproblemas de um nível somam no máximo $n$. O ramo mais curto ($n\to n/3\to\cdots$) tem $\log_3n$ níveis completos e o mais longo ($n\to2n/3\to\cdots$) tem $\log_{3/2}n$ níveis. Logo $n\log_3n\le T(n)\le n\log_{3/2}n+O(n)$: $\Theta(n\log n)$.

\resp{sc-pior8} $W(1)=0$ e $W(n)=2W(n/2)+(n-1)$, pois a intercalação de duas listas de tamanho $n/2$ faz no máximo $n-1$ comparações. $W(2)=1$, $W(4)=5$ e $W(8)=17$. Em geral, $W(n)=n\log_2n-n+1$ para $n$ potência de 2.

\resp{sc-max} Todo elemento que não é o máximo precisa ``perder'' pelo menos uma comparação; senão, o algoritmo não teria como descartá-lo, e um adversário poderia torná-lo o maior. Cada comparação tem um único perdedor, logo são necessárias pelo menos $n-1$ comparações. A árvore de decisão dá só $\log_2n$, porque há apenas $n$ respostas possíveis --- um limite fraco.

\resp{sc-iguais} Com todos os elementos iguais, o teste $V[j]<x$ nunca é verdadeiro: o pivô fica na primeira posição e as partes têm tamanhos $0$ e $n-1$. Então $T(n)=T(n-1)+\Theta(n)=\Theta(n^2)$. Correção: particionar em três partes ($<x$, $=x$ e $>x$) e não recursar na parte dos iguais, ou usar uma partição que divide os iguais entre os dois lados.

\resp{sc-insercao} Ao inserir $V[i]$, ele passa por cima exatamente dos elementos anteriores maiores que ele, isto é, das inversões $(j,i)$ com $j<i$. Cada deslocamento desfaz uma inversão e nenhuma é criada. O total é o número de inversões: até $\binom n2$ no pior caso (vetor invertido) e $\binom n2/2$ em média, pois cada par está invertido com probabilidade $1/2$. Logo $\Theta(n^2)$ nos dois casos.

\resp{sc-intercalar} Considere a entrada $a_1<b_1<a_2<b_2<\cdots<a_n<b_n$. Cada par de elementos vizinhos na saída, $(a_i,b_i)$ ou $(b_i,a_{i+1})$, precisa ser comparado: se um deles não fosse, seria possível trocar os dois valores sem alterar nenhuma resposta das comparações feitas, e a saída ficaria errada. São $2n-1$ pares vizinhos, todos distintos.

\resp{sc-mtrace} As intercalações acontecem nesta ordem: $[27,38]$, $[3,43]$, $[3,27,38,43]$, $[9,82]$, $[9,10,82]$ e, por fim, $[3,9,10,27,38,43,82]$.

\resp{sc-lomuto} Resultado $[2,1,3,4,7,5,6,8]$, e a posição devolvida é 4 (o pivô 4 ocupa sua posição definitiva).

\resp{sc-quickpior} Num vetor ordenado, o pivô (o último elemento) é sempre o máximo do intervalo: uma partição de $k$ elementos faz $k-1$ comparações e deixa uma parte com $k-1$ elementos e outra vazia. O total é $\sum_{k=2}^n(k-1)=\frac{n(n-1)}2$.

\resp{sc-subst} Para $n=1$, $T(1)=0=1\cdot\log_21$. Supondo $T(n/2)=\frac n2\log_2\frac n2$, temos $T(n)=2\cdot\frac n2(\log_2n-1)+n=n\log_2n-n+n=n\log_2n$.

\resp{sc-selecao} A ordenação por seleção faz sempre $\sum_{i=1}^{n-1}(n-i)=\frac{n(n-1)}2$ comparações e no máximo $n-1$ trocas. Não é estável: em $[2_a,2_b,1]$, a primeira troca leva $2_a$ para o fim, depois de $2_b$. A ordenação por inserção é estável, pois um elemento só passa por cima dos estritamente maiores.

\resp{sc-cinco} A árvore de decisão precisa de pelo menos $5!=120$ folhas, logo de altura $\ge\ceil{\log_2120}=7$. Para 12 elementos, $\ceil{\log_2(12!)}=\ceil{28{,}8}=29$.

\resp{sc-kdist} Se cada um dos $k$ valores aparece $n/k$ vezes, há $\frac{n!}{((n/k)!)^k}$ saídas distintas possíveis, e a altura da árvore de decisão é pelo menos o logaritmo desse número, que é $n\log n-n\log(n/k)-O(n)=\Omega(n\log k)$ (usando $\log m!=m\log m-O(m)$). Algoritmo: insira os elementos numa árvore AVL de valores \emph{distintos}, com um contador por valor; a árvore tem no máximo $k$ nós, então cada inserção custa $O(\log k)$; depois, percorra-a em ordem, escrevendo cada valor tantas vezes quanto seu contador. Total $O(n\log k)$.

\resp{sc-nlog2} O nível $j$ tem $2^j$ subproblemas de tamanho $n/2^j$, com custo total $n\log(n/2^j)=n(\log n-j)$. Somando de $j=0$ a $\log n-1$: $n\sum_{j}(\log n-j)=n\cdot\frac{\log n(\log n+1)}2=\Theta(n\log^2n)$. O Teorema Mestre não se aplica porque $f(n)=n\log n$ é maior que $n^E=n$, mas não por um fator polinomial $n^\varepsilon$: cai no ``vão'' entre os casos 2 e 3.

\resp{sc-lineares} (a) $T(n)=1+2+\cdots+n=\frac{n(n+1)}2$. (b) $T(n)=2^n-1$: por indução, $2(2^{n-1}-1)+1=2^n-1$. (c) $T(n)=1+\sum_{i=2}^n\log i=1+\log(n!)=\Theta(n\log n)$.
